\subsection{System and deployment}

We selected MongoDB because it exposes the three consistency controls under
study and provides replica-set failover in one system
\cite{mongodb-replication,mongodb-causal-consistency}. We use three
data-bearing members, initially one primary and two secondaries. This gives
the majority side two copies during an election and lets us compare it with an
isolated former primary without adding an arbiter. Docker Compose fixes the
deployment for repeated trials.

We separate client traffic from replication traffic. This matters in the
partition experiment: we can interrupt replication while leaving the client
able to reach the isolated member. A runner issues the operations; a fault
controller changes member availability or replication reachability. Before
each topology-dependent step, the controller checks which replica-set member
is primary. A post-recovery observer examines all three copies. The
runner has no Docker socket, so fault control is kept outside the client
process.

\maybefigure[fig:architecture]{submission/figures/architecture.pdf}{Client and fault-controller paths for the three-member replica set. During F3, the client can reach the former primary while replication to the other two members is blocked.}

We installed the pinned Python dependencies and used \texttt{make setup} to
start the Compose services, initialise the replica set, and verify its roles
before running a campaign. The setup record supplies the observed versions in
Table~\ref{tab:environment-versions}. Appendix~\ref{app:install} gives the
commands; Appendix~\ref{app:configuration} retains network and image details.

\begin{table}[ht]
\centering
\caption{Software versions recorded for the experiment.}
\label{tab:environment-versions}
\begin{tabular}{@{}ll@{}}
\toprule
Component & Recorded value \\
\midrule
MongoDB image & \texttt{\MongoDBImage} \\
MongoDB server & \texttt{\MongoDBServerVersion} \\
PyMongo & \texttt{\PyMongoVersion} \\
Runner Python & \texttt{\RunnerPythonVersion} \\
Docker Engine & \texttt{\DockerEngineVersion} \\
Docker Compose & \texttt{\DockerComposeVersion} \\
\bottomrule
\end{tabular}
\end{table}
